\section{Ba bài lẽ ra đã khép được khoảng trống}
\label{sec:ch18-ba-bai-le-ra-da-dong}

\index{khoảng trống nghiên cứu!ba bài ngoài danh mục|(}%
Mục trước khép lại ở chỗ chỉ một bài trong chín gọi tên được dụng cụ còn thiếu.
Mục này hỏi câu hiển nhiên tiếp theo: có ai ngoài chín bài ấy đã dựng nó chưa.
Ba tài liệu trả lời được câu ấy, và cả ba đều nằm ngoài danh mục 32 bài của sách.

\index{khoảng trống nghiên cứu!vì sao phải đọc cả ba}%
Vì sao sách phải mở cả ba thì nói trước, vì quy tắc đọc ngoài danh mục hẹp và ba
bài trong một chương là nhiều. Quy tắc, đặt ra ở
chương~\ref{ch:nut-that-khai-niem} và siết lại ở
chương~\ref{ch:do-trung-thuc-chuoi-suy-luan}, là sách mở một bài ngoài danh mục
khi có một khẳng định \emph{mà chương phải xử}, chứ không phải khi chương chỉ
thuật lại. Khẳng định phải xử ở đây là khẳng định của chính sách, câu mạnh nhất
sách đưa ra: phép đối chiếu với ground truth chưa được chạy ở phía attribution
đặc trưng. Ba tài liệu này là ba chỗ mà một phản ví dụ cho câu ấy sẽ nằm. Không
mở chúng là để nguyên câu ấy ở chỗ nó chưa được kiểm, mà một cuốn sách kết thúc
bằng một khoảng trống thì phải kiểm rằng khoảng trống ấy có thật.

\subsection*{Bài siêu đánh giá phát biểu đúng câu của sách}

\index{siêu đánh giá!bài mà chương 11 hoãn lại}%
Bài thứ nhất là bài mà mục~\ref{sec:ch11-khong-co-tang-day} hoãn quyết định đọc
sang chương này. Bài của chương ấy tự tách mình khỏi nó bằng một vế, rằng đó là
bài toán \tn{siêu đánh giá}{meta-evaluation}, \emph{định lượng chất lượng
attribution}, còn khung của bài thì nhắm vào
\tn{tương tác bậc hai}{second-order interaction}~\autocite{p23metagame}. Một bài định lượng chất lượng attribution là đúng thứ
mà khoảng trống này nói là chưa có.

Đọc nó thì thấy đối tượng đúng như vậy. Bài siêu đánh giá lấy chính các chỉ số
chấm điểm lời giải thích làm đối tượng, gọi chúng là các bộ ước lượng chất lượng,
và trong số các chỉ số nó đem ra xét có những chỉ số thuộc đúng các họ mà
chương~\ref{ch:do-do-trung-thuc} dựng lại, trong đó có một chỉ số xóa dần từng
điểm ảnh~\autocite{mmetaquantus}. Vậy nếu có tài liệu nào đã chạy phép đối chiếu
mà sách nói chưa ai chạy thì phải là tài liệu này.

\index{siêu đánh giá!chỗ bài tự nói nó không làm được}%
Bài nói ngược lại, và nói ngay ở chỗ nó dựng khung. Mục 2.2 của nó mang nhan đề
\emph{thách thức của tính không kiểm chứng được}, và đoạn quyết định nằm ở trang
5~\autocite{mmetaquantus}:

\begin{quote}
Một quan sát then chốt khác là ta không thể xác định độ chính xác hay tính hợp lệ
của một bộ ước lượng, tức là nó có thật sự đo cái tính chất mà nó nhắm tới hay
không, bởi một phép đánh giá như vậy đòi phải truy cập được các nhãn ground
truth. Tuy nhiên, vì phân tích độ tin cậy không phụ thuộc vào việc có sẵn nhãn
ground truth, ta vẫn có thể nghiên cứu độ tin cậy của một bộ ước lượng, tức mức
nhất quán chung của nó.%
\footnote{Nguyên văn: \enquote{Another key observation is that we cannot determine
the accuracy or validity of an estimator (i.e., whether it actually measures the
intended quality) since such an assessment requires access to ground truth labels.
However, as reliability analysis does not depend on the availability of ground
truth labels, it is still possible to study the reliability of an estimator, which
refers to its overall consistency.}}
\end{quote}

Câu ấy nói bài không làm được gì; câu tiếp theo nói bài sẽ làm gì thay vào đó, và
nó mở đúng mục mà bài dựng khung của mình, ở trang
sau~\autocite{mmetaquantus}:

\begin{quote}
Không có thông tin ground truth thì ta không thể kiểm định hay tối ưu các bộ ước
lượng theo cái ta muốn chúng đạt được, nhưng ta có thể thay vào đó phát biểu
những tình huống biên hoặc những hành vi mà ta không muốn chúng có.%
\footnote{Nguyên văn: \enquote{Without ground truth information, we cannot validate
or optimise the quality estimators against what we want them to fulfil, but we can
instead articulate edge-case scenarios or behaviours that we do not want them to
exhibit.}}
\end{quote}

Hai câu ấy là câu của mục~\ref{sec:ch09-doi-tuong-khac-nhau}, phát biểu bởi chính
lĩnh vực, trong một bài đã qua bình duyệt, ba năm trước khi bài neo của sách xuất
hiện. Sách đi tìm một phản ví dụ và tìm thấy một nhân chứng.

\index{siêu đánh giá!nước đi giống nước đi của các chỉ số}%
Còn một điều nữa trong bài ấy, và đây là cách đọc của sách chứ không phải câu của
bài. Gặp bức tường vừa dẫn, bài không dừng lại: nó chuyển sang phát biểu các
\emph{chế độ hỏng}, tức những hành vi mà một bộ ước lượng tốt không được
có~\autocite{mmetaquantus}. Nước đi ấy đúng bằng nước đi mà
mục~\ref{sec:ch08-dieu-kien-can} mô tả, ở đó lĩnh vực không kiểm được rằng một
bản đồ nổi bật đúng nên chuyển sang kiểm những điều kiện cần mà nó phải thỏa. Chỗ
giống nhau không phải trùng hợp: hai bên gặp cùng một bức tường, không có ground
truth, và cùng đi vòng bằng cách đổi điều kiện đủ lấy điều kiện cần. Bài siêu
đánh giá làm việc ấy với các chỉ số, còn các chỉ số đã làm việc ấy với các lời
giải thích. Không bài nào nói ra chỗ lặp lại này, vì mỗi bài chỉ đứng ở một tầng.

\subsection*{Bài khảo sát mà ba bài cùng dẫn là một bản kiểm kê}

\index{khoảng trống nghiên cứu!bài được ba bài khác cùng dẫn}%
Bài thứ hai là bài mà ba trong bốn bài lõi của Phần~IV cùng dẫn, mỗi bài cho một việc
khác nhau, và mục ghi chép của sách đã hoãn quyết định đọc nó sang chương này.
Một tài liệu được ba bên độc lập dẫn tới là chỗ đáng hỏi, nhất là khi nó là một
bài tổng quan về \emph{cách lĩnh vực đánh giá lời giải thích}.

Đọc nó thì câu trả lời ngắn: đó là một bài tổng quan hệ thống, và nó không chạy
thí nghiệm nào. Nó gạn từ 606 bài xuống 494 rồi xuống 361 bài được đưa vào, trong
đó 312 bài giới thiệu một phương pháp XAI mới~\autocite{mnauta}. Việc nó làm là
phân loại và đếm những gì người khác làm.

\index{khoảng trống nghiên cứu!Correctness là của lời giải thích, không phải của chỉ số}%
\code{Correctness} là chữ mà bài dùng cho tính chất tương ứng với độ trung thực,
và đó cũng đúng là chữ mà khảo sát của chương~\ref{ch:lime-nao-dang-tin} dùng cho
tính chất đầu tiên trong bảng đánh giá của nó. Bài định nghĩa chữ ấy là mức mà
lời giải thích trung thực với mô hình dự đoán mà nó giải
thích~\autocite{mnauta}, tức là định nghĩa~\ref{def:ch01-do-trung-thuc} nói bằng
lời khác. Khi mô tả họ xóa dần, bài
viết rằng trong một lời giải thích đúng đắn thì điểm số quan trọng của đặc trưng
phải tỉ lệ với độ dịch chuyển của phân bố đầu ra~\autocite{mnauta}. Tiêu chí ở
đó là chính đầu ra của mô hình, tức là \tn{đại lượng thay thế}{proxy} mà cả
chương~\ref{ch:do-do-trung-thuc} dành để mổ, và bài mô tả nó mà không hỏi nó có
đứng vững không.

Câu hướng tới tương lai của bài thì nằm ở phần bàn luận, và tân ngữ của nó là chỗ
quan trọng: bài viết rằng các \emph{phương pháp} XAI cần được kiểm định rộng rãi
để bảo đảm chúng đáng tin và hữu ích~\autocite{mnauta}. Thứ được đem đi kiểm định
là phương pháp giải thích, còn thứ dùng để kiểm định là chính bộ công cụ mà bài
vừa kiểm kê. Đó là một tầng, và tầng trên nó, câu hỏi bộ công cụ ấy đã được kiểm
định bằng gì, bài không đặt ra.

Một con số của bài đáng giữ và một cách đọc nó đáng cấm. Con số là 312 trên 361.
Cách đọc bị cấm là đọc nó thành một tỉ lệ kiểm định: 361 là số bài trong tập hợp
riêng của bài tổng quan này, gồm những bài có liên quan tới việc đánh giá, nên
312 không nói gì về tỉ lệ trong toàn lĩnh vực. Con số đếm việc kiểm chứng với
người thì thuộc về mục~\ref{sec:ch12-mot-cuoc-kiem-dem} và đã ở đó; hai con số
không so với nhau được.

\subsection*{Bài chuẩn đã có sẵn số liệu, và đọc nó thành chuyện khác}

\index{khoảng trống nghiên cứu!bài tới qua danh mục tài liệu}%
Bài thứ ba không được bài nào trong danh mục 32 bài nhắc tới. Nó tới qua danh mục
tài liệu của bài thứ hai, và nó là bài mà hai quyết định đọc kia sẽ bỏ sót. Nó
công bố một bộ dữ liệu tổng hợp và một thư viện để chấm điểm các phương pháp
attribution đặc trưng, trên những phân bố mà kỳ vọng có điều kiện tính được chính
xác chứ không phải xấp xỉ~\autocite{mxaibench}.

Điều làm nó thành bài đáng đọc nhất trong ba bài là danh sách chỉ số của nó. Nó
chấm bằng năm chỉ số và gọi chúng bằng tên riêng của từng chỉ số:
\code{faithfulness}, \code{monotonicity}, ROAR, \code{GT-Shapley} và infidelity.
Hai trong năm là chỉ số mà chương~\ref{ch:do-do-trung-thuc} đã đi qua và đã thêm
khóa trích dẫn, ROAR và infidelity. Chỉ số thứ tư thì khác hẳn bốn chỉ số kia:
\code{GT-Shapley} là hệ số tương quan Pearson giữa trọng số mà phương pháp gán
với \tn{giá trị Shapley}{Shapley value} chính xác, mà bộ dữ liệu tổng hợp cho
phép tính không sai số~\autocite{mxaibench}. Nghĩa là trong cùng một bảng, trên
cùng những lần chạy, bài đặt bốn chỉ số thay thế cạnh một cột có tham chiếu tới
một đáp án tính được.

\index{khoảng trống nghiên cứu!một cột đối chiếu đặt cạnh bốn cột thay thế}%
Bảng ấy in ở trang 7, và nó có bảy cột: sáu phương pháp đã công bố, cộng một cột
ứng với một phương pháp sinh trọng số ngẫu nhiên, đặt ở đó làm mốc. Hãy đọc hai
hàng của nó. Ở hàng \code{monotonicity}, cả bảy cột nằm gọn giữa 0.517 và 0.562,
và cột mốc ngẫu nhiên đứng ở 0.525, tức là nó không đứng cuối: có một phương pháp
đã công bố xếp dưới nó~\autocite{mxaibench}. Ở hàng \code{GT-Shapley}, cũng bảy
cột ấy trải từ $-0.127$ tới 0.930, và cột mốc đứng ở 0.004, với bốn phương pháp
đứng trên nó và hai phương pháp đứng dưới~\autocite{mxaibench}. Cùng một tập
phương pháp và cùng những lần chạy cho ra hai bức tranh không giống nhau.

Chi tiết sắc nhất thì chính bài nêu ra và nêu bằng một câu: phương pháp bị chấm
tệ nhất ở hàng \code{GT-Shapley} lại là phương pháp được chấm tốt nhất ở hàng
\code{monotonicity}, tức $-0.127$ ở hàng dưới và 0.562 ở hàng
trên~\autocite{mxaibench}. Bài giải thích luôn vì sao, và lời giải thích của bài
là chỗ mà mục này dừng lại: phương pháp ấy chọn lần lượt các đặc trưng có ảnh
hưởng lớn nhất lên đầu ra, có hoàn lại, mà đó đúng là thao tác mà
\code{monotonicity} kiểm~\autocite{mxaibench}.

\index{khoảng trống nghiên cứu!một chỉ số có thể bị nhắm để thắng}%
Cách đọc sau đây là của sách, không phải của bài. Một chỉ số mà người ta dựng
được một phương pháp để thắng nó bằng chính cấu tạo của phương pháp ấy thì không
còn là bằng chứng về tính chất ghi trên tên nó. Đó là chẩn đoán của
chương~\ref{ch:do-do-trung-thuc} phát biểu lại, và cũng là điều mà bài khảo sát
của chương~\ref{ch:lime-nao-dang-tin} ghi nhận khi nó viết rằng các chỉ số đánh
giá được chọn sao cho khẳng định của bài được xác nhận. Ở đây cả hai hiện ra
trong một hàng đã in.

Vậy tại sao bảng ấy không khép được khoảng trống. Bài đọc nó theo cách khác: mục
mang nhan đề \emph{hiệu năng qua các chỉ số}, và chủ ngữ suốt mục là các phương
pháp giải thích, còn kết luận rút ra là hiệu năng tương đối của chúng thay đổi
mạnh tùy chỉ số~\autocite{mxaibench}. Không chỗ nào trong bài hỏi câu ngược lại,
là các chỉ số ấy có đo đúng thứ chúng ghi trên tên hay không, dù cột có tham
chiếu đã nằm sẵn trong bảng để hỏi câu ấy.

\index{khoảng trống nghiên cứu!vì sao câu của mục 9 chấm 7 còn đứng}%
Còn một lý do nữa, và nó là lý do khiến câu của
mục~\ref{sec:ch09-doi-tuong-khac-nhau} còn đứng sau khi bảng này được in ra.
Đáp án mà cột \code{GT-Shapley} đối chiếu là giá trị Shapley chính xác, tức một
định nghĩa attribution đã được chọn rồi tính không sai số trên một phân bố đã
biết. Nó không phải những thứ mà hành vi của mô hình phụ thuộc vào. Nói cho gọn
thì cột ấy so một lời giải thích với một lời giải thích khác, cái sau được tính
đúng theo một công thức đã chọn; nó không so lời giải thích với mô hình. Muốn đọc
nó thành phép đo độ trung thực thì phải công nhận trước rằng giá trị Shapley là
câu trả lời đúng, mà chương~\ref{ch:shap} và
chương~\ref{ch:attribution-tu-nguyen-ly-dau} dành hai chương để không công nhận
điều đó. Nên bảng ấy giải quyết được câu hỏi phương pháp nào xấp xỉ giá trị
Shapley tốt nhất, và không chạm tới câu hỏi chỉ số nào đo được độ trung thực.

\index{khoảng trống nghiên cứu!ba bài ngoài danh mục|)}%
Ba bài không khép được khoảng trống theo ba cách khác nhau. Bài thứ nhất đứng
đúng chỗ và tuyên bố rằng chỗ ấy
không đứng được, rồi dựng một khung để làm việc mà không cần đứng ở đó. Bài thứ
hai kiểm kê các dụng cụ và hỏi chúng dùng để kiểm cái gì, chứ không hỏi ai kiểm
chúng. Bài thứ ba có sẵn số liệu để hỏi câu ấy và đọc số liệu ấy thành một câu về
các phương pháp. Cộng lại thì khoảng trống không còn là chuyện thiếu số liệu.
